\tikzset{
  ctrlBlock/.style={
    draw,
    rounded corners=2pt,
    align=center,
    minimum height=8mm,
    minimum width=24mm,
    inner xsep=4pt,
    inner ysep=3pt,
    font=\small
  },
  ctrlSmall/.style={
    draw,
    rounded corners=2pt,
    align=center,
    minimum height=7mm,
    minimum width=20mm,
    inner xsep=3pt,
    inner ysep=2pt,
    font=\footnotesize
  },
  ctrlSum/.style={
    draw,
    circle,
    inner sep=0pt,
    minimum size=5.5mm,
    font=\small
  },
  ctrlGroup/.style={
    draw,
    dashed,
    rounded corners=3pt,
    inner sep=6pt
  },
  ctrlArrow/.style={
    -{Latex[length=2.2mm]},
    line width=0.45pt
  },
  ctrlDashArrow/.style={
    -{Latex[length=2.2mm]},
    dashed,
    line width=0.45pt
  }
}


\begin{figure}[htbp]
  \centering
  \resizebox{0.98\textwidth}{!}{%
  \begin{tikzpicture}
    \node[ctrlBlock, minimum width=31mm] (mobile) at (0,0) {移动机器人\\平台需求};
    \node[ctrlBlock, minimum width=31mm] (wheelleg) at (4.0,1.4) {轮足复合\\高效移动/越障};
    \node[ctrlBlock, minimum width=31mm] (reconfig) at (4.0,0) {模块化自重构\\构型适应/冗余};
    \node[ctrlBlock, minimum width=31mm] (multi) at (4.0,-1.4) {多机器人协同\\信息交互/分布式};

    \node[ctrlBlock, minimum width=36mm] (target) at (8.4,0) {链式自重组\\轮足机器人系统};
    \node[ctrlSmall, minimum width=25mm] (single) at (12.2,1.35) {单体平衡\\变腿长控制};
    \node[ctrlSmall, minimum width=25mm] (chain) at (12.2,0) {铰接约束\\链式运动学};
    \node[ctrlSmall, minimum width=25mm] (coop) at (12.2,-1.35) {分层协同\\DMPC/ADRC};

    \draw[ctrlArrow] (mobile) -- (wheelleg);
    \draw[ctrlArrow] (mobile) -- (reconfig);
    \draw[ctrlArrow] (mobile) -- (multi);
    \draw[ctrlArrow] (wheelleg) -- (target);
    \draw[ctrlArrow] (reconfig) -- (target);
    \draw[ctrlArrow] (multi) -- (target);
    \draw[ctrlArrow] (target) -- (single);
    \draw[ctrlArrow] (target) -- (chain);
    \draw[ctrlArrow] (target) -- (coop);

    \node[ctrlGroup, fit=(wheelleg)(reconfig)(multi), label={[font=\footnotesize]above:相关研究方向}] {};
    \node[ctrlGroup, fit=(single)(chain)(coop), label={[font=\footnotesize]above:本文关注问题}] {};
  \end{tikzpicture}
  }%
  \caption{链式自重组轮足机器人相关研究方向与本文问题定位}
  \label{fig:ch1_research_landscape}
\end{figure}
